\chapter{前言}

本书的愿望很简单：用一部前后连贯的地球与空间科学著作，涵盖从小学一年级到本科毕业的全部内容，
为世界各地的读者而写。

本书力求简洁而严谨：课程由定义、例子、命题、方法和图解构成，只要在相应阶段可以理解，
就给出细致的推理，随后是精心挑选的练习。未作完整推导而给出的结果都明确标为“承认”。
所有练习的完整解答集中在书末。书中引用的有关真实地球、海洋、大气和天空的每个数字都来自参考资料，
根据真实测量绘制的每条曲线都由已发表的数据计算得出。

课程共分五册。第一册涵盖小学和初中，第二册涵盖高中三年，后三册对应大学三年。
两册中小学用书绝不重复讲授同一概念：每个概念只讲一次，放在最适合它首次出现的学年；
后面的章节如需用到，会先以简短的提示\emph{你已经知道}加以回顾，再继续深入。
凡属于物理学、生物学或化学的概念，由同一系列的姊妹篇讲授，本书先作回顾，再从地球科学家的视角展开。
在中小学用书中，每个学年构成一个部分，再分为若干章；每册大学用书对应一个学年，分为若干章。
学年越早，面向的读者越年幼：插图更多，演算步骤更详细，练习也更平缓。

\section*{如何阅读本书}

各陈述在每章内编号，并以颜色区分：定义为蓝色，定理为红色，命题及其同类为橙色，
方法（处理常见任务的简短步骤）为绿色。练习按难度分级，从直接应用的（$\star$）
到较难的（$\star\star\star$）。方框中是不属于论证主线的内容：对前面章节所学内容的回顾，
地质学家和其他地球科学家在野外或实验室中的工作方式，一段历史，
逐年变化的数据（并注明数据所属年份），以及安全提示。

\section*{安全}

认识地球要靠观察：天空、天气、风景、建筑的石材、花园的土壤。这就是本书对读者的全部要求。
切勿直视太阳，无论是否借助仪器：那会致盲。采石场、悬崖、洞穴、洪水中的河流和火山
是受过训练、配备装备的专业人员工作的地方，本书描述他们的工作，但不会让任何人前往这些地方。

\section*{参与贡献}

本书的源代码存放在一个公开的 git 仓库中：
\begin{center}
\url{https://github.com/bvirrion/one-earth-and-space-sciences-book}
\end{center}
欢迎各种贡献，包括新的章节、勘误、更好的解释和更多的练习；
详见仓库根目录下的 \texttt{CONTRIBUTING.md}。
